% 03-Apparatus.tex <- v1 sec:methods-data. Numbers verbatim from v1.
% NO claim block: the absence is structural.

\chapter{The atlas, the model and the rules of measurement}
\label{sec:methods-data}
\framedecl{ER}

\begin{question}
What has to be settled before a score on this atlas means anything? Four
things. What is measured, on what population, against what reference, with what
uncertainty. Tahoe answers none by convention, so each is settled below by a
measurement, and each narrows a claim later.
\end{question}

\section{The panel is a budget, not a biology}

Tahoe-100M~\cite{tahoe100m} holds drug-treated single-cell transcriptomes across
cancer cell lines, with plate-matched DMSO controls. Metadata supplies cell line
(Cellosaurus ID), drug, dose (\texttt{drugname\_drugconc}), plate, and a
fine-grained mechanism-of-action annotation (\texttt{moa-fine}), six fields that
are the whole of what any prompt below can carry.

All analysis is restricted to the $G = 946$ genes in the intersection of the
L1000 landmark set~\cite{l1000} with Tahoe's measured genes, library-normalised
to counts per $10{,}000$ and transformed as $\log_{10}(1 + x)$. The restriction
is a budget choice, not a biological one: the context is $8192$ tokens and a
prompt carries an instruction line and two cell sentences, so the transcriptome
does not fit. Normalisation multiplies every gene in a cell by one constant and
the logarithm is monotone, so within a cell the ordering, and therefore the
sentence, is invariant to both.

\paragraph{What it does change} Everything downstream of a single cell. Scaling
by library size improves comparability across cells of different depth,
but the profiles remain compositional, not absolute, so pseudobulk means, the
drug-specific residual, and every cosine and Euclidean quantity built on them
depend on the normalisation even though the sentences do not. Depth also shapes
which genes are detected, though never their order.

\section{Every treatment sits in one well}
\begin{defn}{Condition, treated well, batch unit}
  A \emph{condition} is a (drug, cell line, plate, dose) tuple, and it is the unit
  everything below is scored at. The physical assignment is the \emph{treated
  well}, so a condition sits inside an assignment and is not one. A plate carries wells from many cell lines at once, and a control pseudobulk is
  matched within both, so the batch unit is the pair (cell line, plate), never the
  plate label alone. A \emph{pseudobulk} is a group's mean
  expression profile; a \emph{control pseudobulk} is the same over plate-matched
  DMSO cells.
  \end{defn}

Tahoe does not treat cell lines one at a time. Roughly $49$ of them are pooled
into a single \emph{well} and one drug at one dose is applied to the pool, so
one well yields around $49$ (drug, dose, cell line) conditions at once and
whatever happened to that well happened to all of them.

The conditions this thesis evaluates --- one drug, at one dose, on one cell
line, on one plate, which is the unit every number below is attached to ---
all carry the plate labels \texttt{plate4}, \texttt{plate5} and \texttt{plate6}.
Labels run within cell line rather than globally, so these are three labels and
not three plates. The cache holds further labels, none of them evaluated.
Across it are $289$ treated wells, recovering a mean of $48.6$ cell lines each
(median $49$, range $43$--$50$). Almost none of them repeat: $286$ of the $287$
distinct (drug, dose) treatments occur in exactly one well.

Those numbers describe this cache, not the atlas: Tahoe includes further plates,
among them a replication plate, so ``one well'' below means one well
\emph{here}. The single exception, one treatment spread over three wells, is
also the whole of the cross-plate structure at condition level, and it accounts
for the $49$ (drug, dose, cell line) groups the unit audit finds on more than
one plate. Every other condition on more than one plate differs in dose, and
each drug is assayed on its own plate or plates with plate-matched controls.

\paragraph{Three consequences} Conditions from one well share a physical
assignment, so they are not independent (\cref{sec:apparatus-clusters}); drug
and plate are partially confounded, so discrimination is measured within plate
(\cref{sec:res-blind}); and a (drug, cell line) pair cannot be held out cleanly
(\cref{sec:methods-holdout}). \cref{fig:confound} draws the three in that
order.

\begin{figure}[htbp]
% MARGINALIA (06-figures.md): fits the text column; caption in the margin.
\begin{lrbox}{\tvfb}%
% ===========================================================================
% figs/confound.tex -- fig:confound, "One assignment, many conditions".
% Figure BODY ONLY. \input this inside a figure environment; the caption and
% \label stay in Sections/03-Apparatus.tex.
%
% SCHEMATIC. Nothing here is plotted from data. The three quantities it names
% are the ones the surrounding text states: roughly 49 cell lines per treated
% well, the other roughly 48 that stay in training when one pair is held out,
% and plate4/plate5/plate6. The squares are twelve per well for legibility,
% which the caption says out loud so the count is not read as a datum.
%
% ---- METRIC REWORK, and why the numbers below are what they are ------------
% The v1 drawing measured 354.51pt (124.6mm) inside a 404.03pt (142mm) block:
% it under-filled the measure by 12%. Fixing that is NOT a uniform \scale of
% the picture, and it is not the 1.44 a naive width ratio suggests, because
% the picture is half drawing and half fixed-size type. Only the drawing
% scales; the three right-hand glosses and the (a)/(b)/(c) labels are set in
% \scriptsize and must stay at that size, since they already print at the
% document's apparatus size.
%
% Measured with pdflatex over this preamble (\the\wd of a savebox):
%   widest gloss node  "a cross-plate drug ... plate contrast" = 181.46pt
%   "held out" node half-width                                 =  18.75pt
% With k the scale on the drawn geometry, the bounding box spans
%   minx = 0.42k*28.4527 - 18.75      (the "held out" label overhangs left)
%   maxx = (5.65k + 0.4)*28.4527 + 181.46
%   width = 148.807k + 211.596 pt
% so k = 1.286 lands the picture at 402.9pt. VERIFIED at 402.96pt; the block
% is 404.03pt. A pure similarity transform of the drawing: x and y take the
% same unit so the lanes keep their aspect, and the squares' minimum size is
% scaled with them (3.0mm * 1.286 = 3.86mm) so a widened lane does not end up
% holding a sparse row of undersized marks. Coordinate VALUES are untouched --
% every number below is v1's -- which is what keeps this a metric change and
% not a redraw.
%
% Gloss alignment: all three glosses sit at one x, the rightmost lane's right
% edge (5.65 units) plus a physical 4mm. 4mm is 0.4cm/1.286 = 0.311 units, so
% x = 5.961. Strip (b)'s own right edge is 5.50, short of (a) and (c) at 5.65;
% hanging its gloss off its own edge would step the left margin of the three
% glosses in and out for no reason.
%
% NEVER \resizebox this picture. Changing its printed size means changing k
% here and re-measuring, because the type must not move.
% ===========================================================================
\begin{tikzpicture}[x=1.286cm, y=1.286cm, font=\small,
  lane/.style={draw=rulegrey, line width=0.4pt, rounded corners=1pt},
  cellb/.style={draw=rulegrey, line width=0.3pt, fill=panelbg,
                minimum size=3.86mm, inner sep=0pt},
  outb/.style={draw=accent, line width=0.7pt, fill=white,
               minimum size=3.86mm, inner sep=0pt},
  lbl/.style={font=\scriptsize, text=quietink},
  strong/.style={font=\scriptsize\bfseries, text=ink}]

% (a)
\node[strong, anchor=west] at (0,3.95)
  {(a)\quad one drug at one dose enters one well};
\draw[lane] (0.15,3.10) rectangle (5.65,3.60);
\foreach \i in {0,...,11}{\node[cellb] at (0.42+\i*0.44,3.35) {};}
\node[lbl, anchor=west] at (5.961,3.35)
  {$\approx 49$ conditions, one assignment};

% (b)
\node[strong, anchor=west] at (0,2.55)
  {(b)\quad each drug is assayed on its own plate};
\foreach \p/\x in {4/0.15, 5/2.00, 6/3.85}{
  \draw[lane] (\x,1.60) rectangle (\x+1.65,2.10);
  \node[lbl] at (\x+0.825,1.85) {\texttt{plate\p}};}
\node[lbl, anchor=west] at (5.961,1.85)
  {a cross-plate drug contrast is also a plate contrast};

% (c)
\node[strong, anchor=west] at (0,1.05)
  {(c)\quad a held-out pair leaves its well in training};
\draw[lane] (0.15,0.35) rectangle (5.65,0.85);
\node[outb] at (0.42,0.60) {};
\foreach \i in {1,...,11}{\node[cellb] at (0.42+\i*0.44,0.60) {};}
\draw[accent, line width=0.4pt] (0.42,0.32) -- (0.42,0.10);
\node[lbl, anchor=north, text=accent] at (0.42,0.08) {held out};
\node[lbl, anchor=west] at (5.961,0.60)
  {its other $\approx 48$ conditions stay in};

\end{tikzpicture}%
%
\end{lrbox}
\sidecap[The well is the unit the atlas actually assigns]{figure}{\captiontitle{One
assignment, many conditions.} The roughly $49$ conditions from one well share
an assignment and a batch (a). Each drug is assayed on its own plate or plates,
so a drug contrast across plates is also a plate contrast (b). Removing a (drug,
cell line) pair leaves its well in training through the other $\approx 48$ lines
(c). Squares are conditions, twelve per well for legibility.\label{fig:confound}}

\end{figure}


Dose is a molar concentration resolved from \texttt{drugname\_drugconc}; An inspection of all $289$ treated wells found $289$ usable molar doses
with zero unit collisions, so the defect was a demonstrated possibility that did
not materialise. 


\section{Barely half the atlas is identifiable}

\cref{tab:scale} states the quantities every reported sample size, $n$, in the investigation should be
located against. One symbol in it runs ahead of its
derivation.\seealso{\cref{sec:res-metrics} derives $\NIR$ and audits its chance
level with empirical control arms.} $\NIR$, the normalised inverse rank, is the
fraction of a declared comparison set whose truth a prediction resembles less
than its own, so chance is $0.50$ under the exchangeability condition of
\cref{sec:res-metrics}.

\begin{scopelimit}[carried into \cref{sec:res-blind} and \cref{sec:res-encoding}]
Two populations are used and they are not interchangeable. The broad
characterisation (\cref{sec:res-blind}) streams a wide sample of the atlas, over
more plates and drugs than cached, so the drug and condition counts in its block
of \cref{tab:scale} count that sample, not the cache. The residual arm
(\cref{sec:res-encoding} onward) works from a smaller cached subset. A count
from one block does not bound a count in the other.
\end{scopelimit}

\begin{table}[htbp]
% MARGINALIA: blocks A and B are key-value lists with dot leaders, set two-up;
% C and D stay tabular. One hairline under each block head; the five lettered
% notes and the general note set below, in the sans.
\caption{\captiontitle{The populations behind every $n$ quoted in the investigation.}
A is a streamed sample of the atlas, B the cache built once and reused. C and D
are two different evaluation schemes, keyed at different units; neither is a
subset of the other.}
\label{tab:scale}
\begin{fullwidth}
% minipage: gives the two-up blocks the full 182mm \hsize inside threeparttable
\begin{minipage}{\linewidth}
\begin{threeparttable}
\newcommand{\blockhead}[2]{\par\noindent{\thead{#1}\quad{\color{quietink}#2}}\par
  \nobreak\vspace{3pt}\noindent{\color{ink}\rule{\linewidth}{0.4pt}}\par\nobreak\vspace{6pt}}
% key-value rows as paragraphs: a long label uses the whole half-width and the
% value sits flush right on its last line, after an unbreakable dot leader.
\newcommand{\kvlead}{\nobreak{\color{rulegrey}\nobreak\cleaders\hbox to 4pt{\hss.\hss}\hskip 0.6em plus 1fill}\nobreak\hskip 0.25em\nobreak}
\newcommand{\kv}[3][0pt]{\par\noindent\leftskip=#1\relax\hangindent=0.6em\relax#2\kvlead\mbox{#3}\par}
\newcommand{\kvblock}{\RaggedRight\parskip=3pt\relax\fontsize{10.5}{13}\selectfont}
\noindent\begin{minipage}[t]{\dimexpr0.5\linewidth-6mm\relax}
\blockhead{A. Task characterisation}{streamed sample of the atlas}
\kvblock
\kv{conditions characterised}{\num{6628}}
\kv[1em]{statistically active against DMSO\tnote{a}}{3,647 / 5,014 (72.7\%)}
\kv[1em]{identifiable at the tested pseudobulk budget\tnote{b}}{3,064 (46.2\%)}
\kv[1em]{with a plate-mate closer than their own split-half}{78.7\%}
\vspace{10pt}
\kv{median cells per condition}{44}
\kv{median \#DEGs per condition (BH $q<0.05$)}{0}
\kv{median SNR (effect $\div$ within-well split-half noise)}{0.75}
\vspace{10pt}
\kv{drugs with a real dose series}{218 / 351 (62.1\%)}
\kv{drugs seen in $\geq 3$ cell lines}{227}
\end{minipage}\hfill
\begin{minipage}[t]{\dimexpr0.5\linewidth-1mm\relax}
\blockhead{B. Residual-frame cache}{rebuilt targets, \cref{sec:methods-residual}}
\kvblock
\kv{cells, treated / control}{620,564 \ (600,000 / 20,564)}
\kv{drugs / cell lines}{106 / 50}
\kv{(cell line, plate) batch groups behind them}{119}
\vspace{10pt}
\kv{conditions with $\geq 40$ treated and $\geq 20$ control cells}{6,617}
\kv[1em]{retained by the data-resolved reliability line}{2,523 of 4,999 train (50.5\%)}
\kv{training / well-held-out validation examples}{147,710 / 3,600}
\end{minipage}
\par\vspace{10pt}
\blockhead{C. Evaluation, generation and expression frames}{keyed at the (drug, cell line) pair}
\noindent\begin{tabularx}{\linewidth}{@{}Y r r r@{}}
 & \thead{Tier 1} & \thead{Tier 2} & \thead{Tier 3} \\
% MARGINALIA: the block-head hairline is this block's one rule; no \midrule.
\addlinespace[4pt]
withheld from training\tnote{c} & nothing & the drug & the pairing \\
conditions available & 4,325 & 2,188 & 99 \\
scored, generation evaluations & \multicolumn{3}{r}{400 per tier} \\
scored, expression-frame $\NIR$, within plate & & 606 & \\
\end{tabularx}
\par\vspace{10pt}
\blockhead{D. Evaluation, residual frame}{keyed at the treated well, \cref{sec:methods-holdout}}
\noindent\begin{tabularx}{\linewidth}{@{}Y r r r@{}}
 & \arm{train} & \arm{unseen_combo} & \arm{unseen_drug} \\
\addlinespace[4pt]
withheld from training & nothing & whole wells & all its wells \\
conditions after inventory filters & 4,999 & 900 & 711 \\
conditions scored, \num{1394} in all & 300 & \needsaudit{895} & 199 \\
treated wells behind them\tnote{d} & 136 & 36 & 28 \\
drugs, 90 distinct in all & 76 & 34 & 10 \\
cell lines, 42 distinct in all\tnote{e} & 39 & 42 & 40 \\
\end{tabularx}
\begin{tablenotes}[flushleft]
\item[a] Over the $5{,}014$ conditions whose (cell line, plate) group carries
plate-matched DMSO cells, not all $6{,}628$.
\item[b] Own within-plate split-half $\NIR \geq 0.8$; per-condition, and distinct
from the aggregate ceilings of \cref{tab:ceilings}.
\item[c] Tier 2 is DrEval's Leave-Drugs-Out. Tier 3 is Leave-Pairs-Out, which
\cref{sec:methods-holdout} declines as an estimand: withholding a pairing leaves
that pair's own well in training through its other $\approx 48$ cell lines.
Block D exists because of it.
\item[d] No well contributes conditions to two splits.
\item[e] Nearly every cell line appears in all three splits, by design: D
withholds wells, not lines, so it asks about an unseen treatment well of a seen
drug and never about transfer to an unseen line.
\item The \texttt{unseen\_combo} count is flagged pending a cluster rerun; no
argument turns on its exact value.
\end{tablenotes}
\end{threeparttable}
\end{minipage}
\end{fullwidth}
\end{table}

\paragraph{The upper block is a result} It describes Tahoe-100M, not this
project. Across the $5{,}014$ control-reachable conditions the median has no
differentially expressed gene at BH $q<0.05$ and a signal-to-noise ratio of
$0.75$ against its own within-well split-half noise, while $72.7\%$ are active
on the whole profile, leaving $27.3\%$ statistically inert. A per-gene test is
far harsher, and $27.3\%$ alone would make the atlas sound better resolved than
it is. The typical condition carries a real response
measured with too few cells, and not no response (\cref{sec:res-blind}).

The identifiability bar is per-condition: a condition's own split-half $\NIR$
against its plate-mates has to clear $0.8$. Of the $6{,}628$ conditions,
$46.2\%$ clear it, and $78.7\%$ have a plate-mate lying closer to them than
their own second half does. Identifiability is a property of the pairing, not of
the drug: over the $227$ drugs assayed in at least three cell lines, the median
drug's split-half reference spans $0.83$ between its most and least identifiable
context. A benchmark omitting those figures
quotes accuracies over a population barely half of which is empirically
identifiable at the tested pseudobulk budget, and a per-drug leaderboard on it
ranks pairings.

% The term is first used in tab:scale above, inside a float, where a marginpar
% cannot be set. This is its first prose use and the chapter's definition point,
% per the \gloss convention (definition at first use in each chapter). Wording
% is the one 05-Limitations.tex already uses, so the two agree word for word.

The three tiers of the last block are defined by what training contained.
\emph{Tier 1} holds conditions whose (drug, cell line) pairing was
seen in training, so it measures fit and not generalisation. \emph{Tier 2} holds
conditions whose drug never appeared, Leave-Drugs-Out in the vocabulary of
\cref{sec:lit-dreval}. \emph{Tier 3} holds pairings whose drug and cell line
were each seen but never together, Leave-Pairs-Out; it is small ($99$
conditions) and is reported only where adequately powered. The residual arm
builds its own, larger holdout instead: \cref{sec:apparatus-vocab} names its
three splits, and \cref{sec:methods-holdout} settles the unit they are keyed
at.

\section{The backbone has seen no perturbation}

A cell is written as a \emph{cell sentence}, its expressed panel genes in
descending order, terminated by \ENDCELL{}, which is registered as an atomic
special token. The prompt supplies cell line, drug, dose, mechanism and the
control cell's sentence; what the model must produce is the treated cell's. 

The terminator is not inherited. A cell expresses only part of the panel, so
sentences differ in length, and a model with no way to mark where its expressed
list ends leaves the boundary between what it ranked and what it omitted to be
inferred from the output. An earlier build here removed the problem rather than
solving it, writing all $946$ panel genes with the unexpressed appended in a
fixed tail: nothing was ever absent, so every convention for handling absence
collapsed into the same one, and most of the generation budget went on a tail
that carried no information. This build lists only what is expressed and closes
it with \ENDCELL{}, an atomic special token registered in the trainer, so a
sentence is variable in length and its end is declared rather than deduced.
Stripping the tail without a marker would have left the same inference to be made
silently. Absence is now a real property of a prediction, which is also why the
rank-based metrics of \cref{sec:res-metrics} need a stated convention for it.

The backbone is \texttt{C2S-Scale-Pythia-1b-pt}~\cite{c2sscale}, a C2S-Scale
checkpoint on the Pythia-1b decoder~\cite{pythia}, cold-started for every arm
from the same base checkpoint with identical hyperparameters (1 epoch, learning
rate $10^{-5}$, gradient accumulation 16, maximum length 8192,
bf16).\seealso{\cref{sec:appendix} tabulates the training arms.} Cold-starting
makes the arms comparable; the training target or objective is the only
variable.

\paragraph{What the checkpoint has seen} It is pretrained on over 50 million
human and mouse cells from 825 single-cell datasets in CELLxGENE and the Human
Cell Atlas, an atlas corpus with no chemical or genetic perturbation data. Its
model card lists cell and tissue prediction, conditioned generation and gene-set
conversion, but not perturbation prediction. The base model arrives at
fine-tuning here having never seen a drug-treated cell.

\paragraph{Two channels stay open} The decoder underneath is Pythia-1b,
pretrained on text in which drug names, targets and mechanisms occur, and every
prompt supplies the drug's \texttt{moa-fine} annotation. A drug the fine-tuning
never saw is not a drug the system knows nothing about.

\begin{scopelimit}[carried into \cref{sec:res-scope} and \cref{sec:limitations}]
``Held out'' therefore means held out from Tahoe fine-tuning, with mechanism
supplied, and every unseen-drug result below is read under that narrower scope.
\end{scopelimit}

\paragraph{Which leaves the target} Two facts sit together here. The checkpoint
has never seen a drug-treated cell, so whatever it comes to know about
perturbation it learns during fine-tuning and learns from the target it is given.
And the formulation it inherits was written for a different job, producing a
plausible cell rather than producing a drug specific treated cell. Cold-starting
fixes the model, so the target is the only thing left that can be moved, and
whether it should be is the question \cref{sec:apparatus-vocab} takes up.

\section{The training target decides the frame}
\label{sec:apparatus-vocab}

The model and the prompt are now fixed, so the string the model is trained to
produce is the only thing left that varies. Two constructions of it run through
the investigation. They are choices about how the target is built and not about
the model, and each carries with it the object a score is taken on, the
competitors it is taken against, and the function that compares them. That triple
is what is called a \emph{measurement frame} here.

\paragraph{The inherited target} The first is the formulation C2S-Scale defines
for perturbation prediction, kept unchanged. The target is the treated cell's
own sentence, its expressed panel genes in descending order, and nothing is
subtracted from it. It is a general-purpose construction, written for a model
asked to produce a plausible cell rather than to tell two treatments apart, and
it is what makes the comparison of \cref{sec:lit-c2s} possible at all.

\paragraph{The perturbation-specific target} The second is built for this
question, because the first one hides the drug: two different drugs on the same cells give
almost the same target string, making it a fault of the data that the model can't give drug specific predictions instead of the model itself.  \Cref{sec:methods-residual} is the measurement
that motivates replacing it. Two subtractions turn a treated profile into a
drug-specific one.

\begin{defn}{The shift, the generic and the drug-specific residual}
For a condition $c$, a (drug, cell line, plate, dose) tuple with at least $40$
treated and $20$ plate-matched control cells, write $d$ for its drug and
\begin{equation}
 \resid(d,c) \;=\;
 \underbrace{\big(\bar{y}_{\textrm{treated}} - \bar{y}_{\textrm{control}}\big)}_{\shift(d,c)}
 \;-\;
 \underbrace{\frac{1}{|D_c \setminus \{d\}|}
   \sum_{d' \in D_c,\; d' \neq d} \shift(d',c)}_{\textrm{generic}(c)} .
\end{equation}
Both bars are pseudobulks: $\bar{y}_{\textrm{treated}}$ averages the condition's
own treated cells and $\bar{y}_{\textrm{control}}$ its plate-matched controls, so
$\resid(d,c)$ is one vector for the whole condition and not a quantity carried by
any single cell. The control subtraction removes the state those cells share,
leaving the \emph{shift}. The mean over the other drugs removes the \emph{generic}
programme every drug triggers, conventionally read as a shared stress and
cell-cycle response, a reading used here for intuition and nowhere tested. What
survives both subtractions is the \emph{drug-specific residual}, and taking it as
the string the model is trained to emit makes it the \emph{residual target}. A
condition therefore supplies many training examples but only one target: the
prompt varies from one plate-matched control cell to the next, and every one of
them is paired with the same residual string.

\smallskip
$D_c$ is the training drugs sharing $c$'s (cell line, plate) group, each drug
entering once whatever its doses. That grouping is the generic's \emph{scope},
and this build takes it at plate scope; \cref{sec:res-transfer} reports the
cell-line-scoped alternative alongside it. Two properties of $D_c$ are fundamental. It is \emph{split-before-fit}: the holdout is fixed from metadata
before any statistic is estimated, the generic included. And it is
\emph{leave-one-drug-out}: a condition's own drug never enters its own generic,
so the subtraction removes what the other drugs share and cannot remove the drug
being predicted.

\smallskip
The scope is also the limit. What remains is not what makes this drug different
from all drugs. It is what makes it different \emph{from the other training
drugs that happen to share its plate and cell line}.
\end{defn}
\paragraph{Writing the residual as a string} $\resid(d,c)$ is a signed vector
over the panel, and a decoder emits tokens rather than vectors, so the target has
to be written down before it can be trained on. The obvious form is the wrong
one. A single list in descending order buries the most strongly repressed genes
at the tail, among the genes the drug did not move at all, and repression is
exactly what the control subtraction was performed to expose. The sign is the
biology, so the encoding keeps it: the target is two ordered blocks, the most
up-regulated genes first, a \DOWN{} separator, then the most down-regulated
genes, closed by \ENDCELL{}. Up to $100$ genes enter each block, which is the
$200$ target positions \cref{tab:divergence} counts over, and a gene enters a
block only if its residual actually carries that sign, so a condition with few
movers emits a short one. \DOWN{} is registered as an atomic special token
alongside \ENDCELL{}; unregistered it splits into sub-words and the up/down
boundary stops being a clean symbol. What the model is asked to produce is
therefore a ranked signed signature, not a cell.

\paragraph{Reading a generation back, in each frame} What the model emits is a
string, and a similarity needs vectors, so a generation --- and the truth it is
scored against --- has to be read back into one. Three read-backs are used below,
and which one applies is a property of the metric and not of the frame: the frame
names the object scored, a cell profile or a residual, and not the units it is
scored in.

The first leaves the sentence where it is. A gene's value is its position in the
sentence, counting from one, and a panel gene the sentence omits takes the rank
the absent-gene convention assigns (\cref{sec:app-convention}). $\DEdr$ and
$\paneltau$ are computed here and never leave rank space.

The second decodes to a reconstructed expression profile. Write $p_g$ for the
position at which gene $g$ appears. Then
\begin{equation}
 \psi(s)_g \;=\;
 \begin{cases}
   \max\!\big(0,\; a\log_{10} p_g + b\big) & \text{if $g$ appears in $s$},\\[2pt]
   0 & \text{if $g$ does not,}
 \end{cases}
\end{equation}
with $a$ and $b$ a single slope and intercept fitted once for the whole run
(\cref{sec:lit-c2s}). This is what \texttt{nir\_expr} reads, prediction and truth
both decoded, and it is the only read-back that invents a magnitude. Note that
the absent-gene convention does not reach it: $\psi$ fills an omitted gene with
zero whatever the convention says.

The third keeps only sign and order, because a residual has no expression units to
decode into. Write $s$ for a signed signature of either kind, and let $U(s)$ and
$D(s)$ be its up and its down block: for a generation, the two lists the model
emitted, in the order it emitted them; for a continuous residual, its most
positive and its most negative entries under the same cut of $100$. Let $r_g$ be
the position of gene $g$ within whichever block holds it, counting from one. Then
\begin{equation}
 \phi(s)_g \;=\;
 \begin{cases}
   +\,1/\log_{2}(r_g + 1) & \text{if } g \in U(s),\\[2pt]
   -\,1/\log_{2}(r_g + 1) & \text{if } g \in D(s),\\[2pt]
   0 & \text{if $g$ is in neither block.}
 \end{cases}
\end{equation}
The leading gene of each block carries weight one and the weight falls away down
the block, so a gene's rank is all that reaches the vector and the value that
earned it that rank is discarded. The discount is what makes the ordering inside a
block count for something --- agreeing on the most up-regulated gene is worth more
than agreeing on the hundredth, where one rank is barely distinguishable from its
neighbours, and a flat encoding would score those two agreements alike.

All three read position and nothing else. A gene's identity decides \emph{where}
a value lands in the vector, never \emph{what} that value is, and all three fill
with zero or with a convention whatever the string leaves out. They differ in what
they then claim: $\psi$ writes position into expression units through a map fitted
to data, so an expression-frame distance taken on it is a distance between two
reconstructions; $\phi$ writes position into fixed weights that no fitting
produced, so a residual-frame comparison is a rank-weighted overlap of gene sets
and carries no units at all. \Cref{tab:sigma} records which read-back each quoted
figure rests on.
\paragraph{Why the target carries a frame with it} A model trained on the first
target emits a cell profile; a model trained on the second emits a signed
up-and-down signature. The two are not the same kind of object, and neither
scoring path accepts the other's output: the rank-to-expression decode is fitted
for expression ranks and its truths are full pseudobulks, while a residual
signature has no expression units to decode into. So the choice of target
settles the space a prediction lives in, and a score has to be taken in that
space or not at all. Comparing a number from one frame with a number from the
other is the most expensive misreading this thesis allows.

Both frames score by discrimination and not by proximity. A prediction is never
asked how close it came to its own truth. It is set beside a lineup of other
conditions' truths and asked whether it resembles its own more than it resembles
any of theirs. That fraction is the \emph{normalised inverse rank}
$\NIR$,\seealso{\cref{sec:res-metrics} derives $\NIR$ and audits its chance level
with empirical control arms.} and the lineup is the \emph{comparison set} $J$.
Choosing $J$ is a substantive decision rather than a formality: anything every
candidate in $J$ has in common cannot help a prediction win, because it is
equally true of all of them. Narrowing $J$ is how a nuisance cue is switched off.

The two frames narrow it differently. The \Eframe{} \emph{expression frame}
scores the inherited target's output, a predicted cell profile, against the other
drugs \emph{on its own plate}, so every candidate carries the same batch and no
prediction can score by recognising the plate it came from --- a leak
\cref{sec:res-blind} measures rather than assumes. The \Rframe{} \emph{residual
frame} scores the residual target's output against the other drugs \emph{in its
own cell line}, more widely scoped than the generic, so every candidate carries
the same cellular background and the drug is the only thing left to separate
them. Chance sits at $0.50$ in both, under an exchangeability condition
\cref{sec:res-metrics} tests with control arms instead of assuming, and $J$ is
declared at every score below. A score without its comparison set is not a score.

The \emph{similarity} $\sigma$, the function by which a prediction is judged
closer to one truth than to another, is the third thing the target settles, and
it follows neither from the scored object nor from the comparison set on its own:
a negated Euclidean distance on decoded expression in the expression frame, a
cosine on the signed rank-discounted residual in the residual frame, the two
instances \cref{tab:sigma} fixes with their comparison sets. A score without its
similarity is not a score either, and both are to be declared before a number is
read.

% ===========================================================================
% frames.tex -- the two measurement frames, EXPR and RESID, as two cards.
% \input by Sections/03-Apparatus.tex. Compiles against thesis_v2/preamble.tex
% and nothing else; no external data, no \includegraphics.
%
% NO DATA SOURCE. This is a schematic. The only numeral on the cards is the
% chance level 0.50, which is a definition and not a measurement: both frames
% are two-alternative retrievals, so chance is 0.50 by construction. It agrees
% with the prose immediately above this figure in Sections/03-Apparatus.tex
% ("Both put chance at $0.50$"), with Sections/02-Background.tex:244 and with
% Sections/04a-ruler.tex:530 ("Chance is $0.50$ by ..."). Nothing here is read
% from RESULTS_cluster/ because nothing here is a result. The numeric version
% of this argument -- the two frames on their own scales -- is fig:two-scales
% in 04d; Chapter 3 is apparatus and comes before those results exist, so this
% figure carries no scale, no axis and no measured quantity. Sibling
% frames.json records the geometry and that single number.
%
% ---------------------------------------------------------------------------
% THE BUG, AND WHY THE FIX IS NOT THE ONE-LINER IT LOOKS LIKE.
%
% Symptom, both cards: the two-line value on row 2 collided with its own key.
% Measured in the built thesis, ink gap between the key and its value, in
% PostScript points: r1 -0.02, r2 -3.02, r3 +2.12. Only r2 is visible as a
% collision. (r1 is the descender of "predicted" level with the ascender of
% "profile"; different x, they do not touch.)
%
% Diagnosis, as briefed and confirmed: val/.style used anchor=WEST, which
% centres a node's box on its y, so a TWO-LINE value grew upward into the key
% above it while one-line values did not.
%
% The briefed fix was anchor=north west with every value at its key's
% y - 0.30cm. Built, that removes the collision and introduces a worse fault:
% it lowers each value until it sits CLOSER TO THE NEXT KEY THAN TO ITS OWN,
% so every value reads as the heading of the row beneath it. Sweeping the
% offset, all measured inside one PDF by tools/_measure_frames.py and recorded
% in frames.json, gave "within" = key to its own value, "between" = value to
% the next key:
%
%   offset      within r1 / r2 / r3      between: after r1 / after r2
%   0.15          2.21  2.24  3.14            9.60        7.43
%   0.30          6.46  6.49  7.39            5.35        3.18   <- reversed
%   0.40          9.29  9.33 10.23            2.51        0.35
%
% and then the same figure BUILT INSIDE THE THESIS at offset 0.15 measured
% 5.95 / 5.98 / 5.29 within against 5.86 / 3.69 between -- 3.7pt away from the
% harness, and ambiguous. The offset is not a constant of the figure.
%
% ROOT CAUSE, and why this file no longer uses align=left. TikZ's `align'
% option typesets the node in a minipage, and a minipage's first-line height
% follows the AMBIENT \baselineskip. So a value node's ink lands in a place
% that depends on the prose that happens to precede the float: about 2.2pt
% between a \normalsize paragraph and a \small line, and about 3.7pt between a
% test harness and Chapter 3. No choice of offset survives an edit to the
% surrounding text. Calibrating one would have hidden the fault, not fixed it.
%
% THE FIX ACTUALLY USED: every node is a single line anchored BASE WEST, so
% each y coordinate below IS that line's baseline. Nothing is centred, nothing
% is measured from a box edge, no minipage is created, and no ascender or
% descender can move a line. The two-line value is two nodes rather than one
% node with `\\'. Key-to-value leading is therefore identical on all three rows
% and on both cards by construction -- which is what the brief asked for -- and
% it is now also identical from build to build.
%
%   key baselines    3.70   2.72   1.52     (row pitch unchanged)
%   value baselines  3.37   2.39   1.19     (key - 0.33cm)
%   row-2 value, second line  2.06          (previous line - 0.33cm)
%
% ONE step governs everything inside a row: 0.33cm = 9.39pt, one \scriptsize
% line of leading. A key and its value therefore set as two consecutive lines
% of a single block, exactly as rows 1 and 3 of the published figure did, and a
% wrapped value continues on the same rhythm. Row SEPARATION is carried
% entirely by the key pitch, which is untouched (0.98cm, then 1.20cm), giving
% 0.65cm and 0.53cm of clear space before the next key -- roughly twice the
% within-row step, so a value binds to the key above it and not to the one
% below.
%
% Verified by reading the baselines back out of the built PDFs. Chapter 3 in
% main.pdf and all four contexts of figs/_frames_harness.tex report the same
% seven baselines to 0.001cm:
%   3.703  3.375 | 2.727  2.398  2.059 | 1.531  1.203   (tikz y, both cards)
% Note the harness and the thesis are only comparable on BASELINES. Glyph
% bounding boxes reported by a PDF text extractor differ by a constant ~2.1pt
% between the two files because the font subsets differ, so ink-gap numbers
% are only meaningful WITHIN one PDF.
%
% OTHER LAYOUT NOTES
%   * cards are 5.9cm wide with a 0.5cm gutter, 12.3cm overall. NOT rescaled:
%     12.3cm sits inside the 14.2cm body block and leaves the page free to
%     carry a margin note.
%   * the dotted rule and the sentence under it span BOTH cards on purpose.
%     They are the only mark that crosses the gutter; everything else is
%     card-local, because the point of the figure is that the columns do not
%     exchange.
%   * if a value ever needs a third line, add a node at (its predecessor
%     - 0.34cm). Do not reintroduce align/`\\'.
% ===========================================================================
\begin{figure}[htbp]
% MARGINALIA (06-figures.md): no borders on the two cards; the frame names are
% sans tags in accent; the figure fits the text column, so its caption is in
% the margin (\sidecap).
\begin{lrbox}{\tvfb}%
\begin{tikzpicture}[x=1cm, y=1cm, font=\small,
  panel/.style={draw=none},
  % base west on both: every y below is a BASELINE. See the header note --
  % anchor=west centres the box (that was the collision) and align=left makes
  % the node's height depend on the surrounding prose (that was worse).
  key/.style={font=\scriptsize\bfseries, text=quietink, anchor=base west},
  val/.style={font=\scriptsize, text=ink, anchor=base west}]

\draw[panel] (0,0.55) rectangle (5.9,4.6);
\draw[panel] (6.4,0.55) rectangle (12.3,4.6);

\node[anchor=base west] at (0.25,4.18)
  {{\sffamily\fontsize{7.5}{9}\selectfont\bfseries\color{accent}\textls[80]{\ding{110}\,EXPR}}};
\node[anchor=base west] at (6.65,4.18)
  {{\sffamily\fontsize{7.5}{9}\selectfont\bfseries\color{accent}\textls[80]{\ding{115}\,RESID}}};

% key baseline 3.70 / 2.72 / 1.52 ; value baseline = key - 0.33 ; wrap - 0.34.
\node[key] at (0.25,3.70) {what is predicted};
\node[val] at (0.25,3.37) {a cell profile};
\node[key] at (6.65,3.70) {what is predicted};
\node[val] at (6.65,3.37) {the drug-specific residual};

\node[key] at (0.25,2.72) {what it is ranked against};
\node[val] at (0.25,2.39) {other drugs on the};
\node[val] at (0.25,2.06) {same plate};
\node[key] at (6.65,2.72) {what it is ranked against};
\node[val] at (6.65,2.39) {other drugs in the};
\node[val] at (6.65,2.06) {same cell line};

\node[key] at (0.25,1.52) {where chance sits};
\node[val] at (0.25,1.19) {$0.50$};
\node[key] at (6.65,1.52) {where chance sits};
\node[val] at (6.65,1.19) {$0.50$};

\draw[ink, line width=0.4pt, densely dotted] (0,0.25) -- (12.3,0.25);
\node[font=\scriptsize\itshape, text=ink, anchor=base] at (6.15,0.02)
  {the same number in the two columns is not the same measurement};
\end{tikzpicture}%
\end{lrbox}
\sidecap[The two measurement frames]{figure}{\captiontitle{Two frames, one chance level, and
no exchange rate between them.} The prediction, the comparison set and the
meaning of a score change between frames; the chance level does not.\label{fig:frames}}
\end{figure}


\section{One estimator, and not a replicate}

Almost every negative result below is stated as a distance from a
\emph{ceiling}. 

\begin{defn}{Within-well split-half precision reference}
The score obtained when a disjoint cell subset of the same treated well is
substituted for the prediction, recomputed for every stratum scored, so its
value moves with the stratum. It is not a replicate ceiling. What varies across
the rows of \cref{tab:ceilings} is the frame, the comparison set and the cell
budget, never the estimator.
\end{defn}

Two properties limit how the ceilings read. First, they are \emph{split-half}
estimates and not biological replicates, because this cache places $286$ of
$287$ treatments in exactly one well. Second, a split-half ceiling estimates
within-well sampling precision, which is optimistic relative to true biological
reproducibility, and every coverage figure in \cref{sec:res-lookup} inherits
that optimism.

\paragraph{The one exception} One ceiling is not that estimator and not a row in
\cref{tab:ceilings}. The positive control inside the expression-frame $\DRF$
audit of \cref{sec:res-metrics} is a denoised interpolated duplicate, not a raw
split-half, and that section turns on the difference.

\begin{table}[htbp]
% MARGINALIA: full width; caption above at the text width, notes below in the
% sans. Text columns widened so every metric description fits on two lines.
\caption{\captiontitle{One estimator, seven settings.} Every ``ceiling'' here is a
within-well split-half precision reference, the score of a disjoint subset of
the same treated well's cells. They differ in frame, comparison set, cell budget
and split contents.}
\label{tab:ceilings}
\begin{fullwidth}
\begin{threeparttable}
\begin{tabularx}{\linewidth}{@{}r L{61mm} L{48mm} Y@{}}
\thead{Ceiling} & \thead{Frame and metric} & \thead{Comparison set and budget}
  & \thead{Where used} \\
\midrule
$0.76$ & expression, $\DEdr$ & two real cells & metric audit
  (\cref{sec:res-metrics}) \\
\addlinespace
$0.625$ / $0.575$ & expression, $\NIR$ (negated Euclidean on decoded expression)
  & leak audit; cross-plate / within plate
  & plate leak (\cref{sec:res-blind}) \\
\addlinespace
$0.67$--$0.83$ & expression, forced choice & within plate, by cell budget
  & \cref{sec:res-blind} \\
\addlinespace
$0.576$ & expression, $\NIR$ (negated Euclidean on decoded expression)
  & within plate, tier 2 & the blindness verdict \\
\addlinespace
$0.61$ / $0.76$ / $0.85$ & expression, $\NIR$ (negated Euclidean on raw
  pseudobulk, undecoded) & \emph{cross-plate};
  $10$/$40$/$120$ cells & identifiability curve \\
\addlinespace
$0.956$/$0.824$/$0.836$ & residual, $\NIR$ (cosine on the signed rank-discounted
  residual) & within cell line;
  train/combo/drug & \cref{sec:res-lookup} \\
\addlinespace
$0.742$ / $0.7432$ & retrieval, negated Euclidean on undecoded log-expression
  & within condition, split-half
  & \cref{sec:methods-residual} \\
\end{tabularx}
\begin{tablenotes}[flushleft]
\item The $\DRF$ positive control of \cref{sec:res-metrics} is the
investigation's one exception and is not a row here.
\item None of these is a biological replicate ($286$ of $287$ treatments sit in
a single well).
\item In the residual frame the raw reference is not constant across splits (a
spread of $0.132$), but that spread is an artefact of how the splits were
filtered and not a difference in recoverable signal; matched at the training
reliability line the three agree to $0.006$ (\cref{sec:limitations}).
\end{tablenotes}
\end{threeparttable}
\end{fullwidth}
\end{table}

\begin{scopelimit}[carried into \cref{sec:res-lookup} and \cref{sec:limitations}]
A distance to the reference is comparable within a split and not between them,
because the splits are scored on populations selected by different rules. And
because the reference is a sampling-precision estimate and not a biological
replicate, every coverage figure computed against it is a distance to an
optimistic bar by an unknown amount. Where one is quoted, the ordering carries
the claim.
\end{scopelimit}

\section{Two crossed dependencies, not one}
\label{sec:apparatus-clusters}

Two rules apply to every result below, plus one statement of provisionality.
Conditions are not independent, and not independent in two crossed ways at once
(\cref{fig:clustering}): two from one treated well share a physical assignment
and a batch, two from one cell line share the biology. Treating them as
independent gives intervals about $1.5$ times too narrow on the scramble
strata and $1.7$ to $2.5$ times too narrow on the contrasts against the lookup
baselines. Clustering on the assignment unit alone uses the larger cluster
count ($289$ wells against $\approx 50$ lines) and would tighten every interval
while sounding more rigorous. That is the failure the estimator below avoids.

% ===========================================================================
% clustering.tex -- the two crossed groupings, plus the drug the wells nest in.
% \input by Sections/03-Apparatus.tex. Compiles against thesis_v2/preamble.tex
% and nothing else; SCHEMATIC -- no external data, no \includegraphics, and
% deliberately NO number anywhere in the drawing. The lattice is 9 columns of
% 6 dots because that is what reads at this measure, not because the cache
% holds 9 wells or 6 lines; the caption says so in as many words.
%
% WHY IT WAS REDRAWN. The version it replaces sat inline in 03-Apparatus.tex
% at lines 477-510 (v1's is under ../thesis/ and is not touched).
%
%   1. THE DOT WAS NOT ON THE LATTICE. The old bands were centred at y=2.6 and
%      x=5.0, while the lattice rows sat at 0.5+0.6r (2.3, 2.9) and the columns
%      at 0.7+1.1c (4.0, 5.1). Both centre lines fell BETWEEN dots, so the
%      accent "one condition" dot was drawn in the gap. The caption said "the
%      dot where they meet is one condition" and the drawing did not show one.
%      Here every coordinate is derived from the two constants the lattice is
%      generated from (\px, \py) and the highlight is addressed by INDEX
%      (\hc, \hr), so band centres and accent dot cannot drift off a lattice
%      point again.
%
%   2. THE BANDS WERE PADDED BY EYE. Old horizontal band 0.35..8.45 against
%      dots at 0.7..7.3 -- 0.35 short at the left, 1.15 long at the right. Old
%      vertical band 0.12..4.08 against dots at 0.5..4.1 -- 0.38 short at the
%      bottom, 0.02 over at the top. Here each band ends exactly half a lattice
%      pitch past its first and last dot (\BL,\BR,\BB,\BT below), computed, so
%      both are symmetric by construction.
%
%   3. THE DOTS WERE TEXTURE, NOT CONDITIONS. They were rulegrey!70, i.e.
%      #9A9A9A at 70% over white ~= #C0C0C0 -- lighter than the document's
%      quiet ink and far too light for a mark the caption asks the reader to
%      count. They are quietink (#6B6B6B) here, and 1.7pt rather than 1.5pt.
%
%   4. IT PAID FOR FULLWIDTH AND DID NOT USE IT. The old float opened
%      \begin{fullwidth} -- 174mm, and by the preamble's own rule a page
%      carrying one may then carry no margin note -- and drew 148.6mm. This
%      draws inside the 142mm measure (see MEASURED below), so it is an
%      ordinary body-width float and its page keeps the margin column.
%
%   5. THE DRUG WAS IN NO FIGURE AT ALL, and it is the axis that flips three
%      verdicts later. It is drawn in a deliberately different register from
%      the two bands: outline only, no fill, rounded corners, against two solid
%      tint rectangles. A well carries one drug at one dose (sec. "Every
%      treatment sits in one well"), so wells NEST inside drugs -- three whole
%      well-columns per envelope -- while the cell-line band CROSSES all three
%      envelopes and runs past the outer edge of both end ones. Nesting and
%      crossing are then the same picture, told apart by shape, not by colour.
%
% GEOMETRY, in cm, all of it derived from \px and \py. Every clearance here is
% computed from those two constants, not judged by eye:
%   lattice     9 cols x 6 rows, pitch 1.25 x 0.81, dots (0,0)..(10.00,4.05)
%   highlight   col 4 -> x = 5.00, row 4 -> y = 3.24     <- both ON the lattice
%   bands       half-thickness 0.25, so 0.155 of ground to the next dot row
%               (0.405 half-pitch - 0.25) and 0.095 clear of that dot's edge;
%               ends at -0.625/10.625 and -0.405/4.455, i.e. exactly half a
%               pitch past the first and last dot of the row and the column
%   envelopes   3 x 3 columns, pad 0.46 -> -0.46..2.96, 3.29..6.71, 7.04..10.46,
%               so 0.33 between neighbours; vertically -0.585..4.635, giving
%               0.18 of margin above and below the well band each one contains
%   callout     the leader leaves the accent dot at 46.3 deg, clears the
%               nearest lattice dot (6.25,4.05) by 0.354, and crosses the
%               middle envelope's TOP edge at x = 6.32, inside that envelope's
%               span 3.29..6.71 -- so it pierces one boundary and not two
%
% MEASURED, not estimated. figs/_harness_clustering.tex \input s this file
% against the real preamble, captures the float into a box with \caption and
% \label gobbled, and \typeout s \wd. Reading:
%   tikzpicture  385.40pt = 135.45mm wide, 187.44pt = 65.87mm tall
%   \textwidth   404.03pt = 142.00mm
%   fills 95.4% of the measure, 18.63pt = 6.55mm of slack.
% The build's overfull-hbox gate is 0; this figure contributes none. The
% fullwidth version it replaces drew 148.6mm against a 174mm allowance -- it
% overran the body measure, which is why it was in a fullwidth in the first
% place, and cost that page its margin column to do it.
%
% TYPE. \scriptsize on every label, matching fig:confound three pages earlier
% and fig:residual-energy in chapter 4. Deliberately not raised: this figure is
% read alongside fig:confound, and a lone larger label set would read as a
% different apparatus rather than as a corrected one.
% ===========================================================================
\begin{figure}[htbp]
% MARGINALIA (06-figures.md): fits the text column; caption in the margin.
\begin{lrbox}{\tvfb}%
\begin{tikzpicture}[x=1cm, y=1cm, font=\small]
  % --- the constants everything else is derived from -----------------------
  \def\px{1.25}   % column pitch  (one well per column)
  \def\py{0.81}   % row pitch     (one cell line per row)
  \def\hc{4}      % highlighted column index, 0..8
  \def\hr{4}      % highlighted row index,    0..5
  \def\bt{0.25}   % half band thickness
  \def\ep{0.46}   % envelope padding around its three columns
  \def\vm{0.18}   % envelope margin above and below the well band

  \tikzset{
    cond/.style={fill=quietink},
    band/.style={fill=rulegrey!30},
    drugenv/.style={draw=quietink, line width=0.5pt, rounded corners=4pt},
    lead/.style={draw=quietink, line width=0.4pt},
    lbl/.style={font=\scriptsize, text=quietink},
    strong/.style={font=\scriptsize\bfseries, text=accent}}

  % centre lines of the two bands: ON the lattice, by construction
  \pgfmathsetmacro{\HX}{\hc*\px}            %  5.000
  \pgfmathsetmacro{\HY}{\hr*\py}            %  3.240
  % band ends: half a pitch past the first and last dot, both sides, both bands
  \pgfmathsetmacro{\BL}{-\px/2}             % -0.625
  \pgfmathsetmacro{\BR}{8*\px+\px/2}        % 10.625
  \pgfmathsetmacro{\BB}{-\py/2}             % -0.405
  \pgfmathsetmacro{\BT}{5*\py+\py/2}        %  4.455
  \pgfmathsetmacro{\EB}{\BB-\vm}            % -0.585
  \pgfmathsetmacro{\ET}{\BT+\vm}            %  4.635

  % ---- 1. the two crossed groupings, as solid tint bands ------------------
  \fill[band] (\BL,\HY-\bt) rectangle (\BR,\HY+\bt);      % one cell line
  \fill[band] (\HX-\bt,\BB) rectangle (\HX+\bt,\BT);      % one well

  % ---- 2. the drug, as an outline envelope over three whole well-columns --
  % It encloses the well band top and bottom as well as left and right, so the
  % well reads as INSIDE the drug and not merely beside it.
  \foreach \g in {0,1,2}{
    \pgfmathsetmacro{\GL}{3*\g*\px-\ep}
    \pgfmathsetmacro{\GR}{(3*\g+2)*\px+\ep}
    \draw[drugenv] (\GL,\EB) rectangle (\GR,\ET);}

  % ---- 3. the conditions --------------------------------------------------
  \foreach \c in {0,...,8}{
    \foreach \r in {0,...,5}{
      \fill[cond] (\c*\px,\r*\py) circle (1.7pt);}}

  % ---- 4. the intersection ------------------------------------------------
  \fill[accent] (\HX,\HY) circle (2.4pt);
  \draw[accent, line width=0.4pt]
    (\HX+0.11,\HY+0.13) -- (\HX+1.66,\HY+1.75);
  \node[strong, anchor=west] at (\HX+1.78,\HY+1.81)
    {one condition: the intersection};

  % ---- 5. direct labels, one per grouping. No legend, no key. -------------
  % cell line: to the left of its own band.
  \node[lbl, anchor=east] at (\BL-0.15,\HY) {one cell line};

  % well: below its own band, on a lead that PIERCES the envelope's lower
  % edge -- the crossing of that boundary is the nesting, drawn.
  \draw[lead] (\HX,\BB) -- (\HX,\BB-0.34);
  \node[lbl, anchor=north] at (\HX,\BB-0.40) {one well};

  % drug: above the leftmost envelope, centred on it. Placed there and not on
  % the middle envelope because the middle one is where the well band and the
  % callout leader already are, and "one drug" centred over the highlighted
  % column would compete with "one well" for the same column.
  \node[lbl, anchor=south] at (\px,\ET+0.06) {one drug};

  % NO free-floating note line. The version this replaces carried "neither
  % grouping is nested in the other" as a caption-width sentence under the
  % lattice. Two reasons it is gone: the caption now states it in the same
  % words, and at page scale a \scriptsize sentence sitting 4mm above
  % "Figure 3.3." reads as the caption's first line rather than as part of the
  % drawing. The claim is not dropped -- it moved into the caption, where the
  % new drug clause needs it anyway.
\end{tikzpicture}%
\end{lrbox}
\sidecap[The two crossed groupings, and the drug the wells nest in]{figure}{\captiontitle{A
condition belongs to a well and to a cell line at once, and the two groupings
cross.} Each dot is a condition; the horizontal band is one cell line, the
vertical band one well. The dot where they meet is one condition, which is why
the intersection term is the ordinary independent-sampling variance. The
rounded envelopes are drugs: a well carries one drug at one dose, so wells nest
inside drugs, while the cell-line band crosses every drug and neither of the
two bands is nested in the other. The lattice is schematic --- nine columns of
six, for legibility, and not a count of anything in the cache.\label{fig:clustering}}

\end{figure}


Every interval on a \emph{mean} is therefore two-way cluster-robust
(Cameron--Gelbach--Miller),
\begin{equation}
 V \;=\; V_{\textrm{line}} \;+\; V_{\textrm{well}} \;-\; V_{\textrm{intersection}},
\end{equation}
where the intersection of the two groupings is the single condition, so
$V_{\textrm{intersection}}$ is the ordinary independent-sampling variance. When
it goes non-positive in a finite sample, the wider one-way variance is used and
the branch recorded. The formula above is valid only where the two groupings
genuinely cross, which cell line and well do: a treated well spans
$\approx 49$ cell lines and a cell line recurs across wells. Drug and well do
\emph{not} cross, because one drug is assigned per well and wells therefore nest
strictly inside drugs. Where an interval is reported on the drug axis the
estimator tests the two groupings for nesting first and, finding it, clusters
one-way on the drug rather than applying the expression above to a nested pair;
the branch taken is recorded beside the interval. Reading ``two-way'' as
``the equation above was applied to drug and well'' would therefore be wrong.
It is asymptotic in clusters, so the reference distribution
is Student's $t$ with one fewer degree of freedom than the smaller cluster
count. The \texttt{unseen\_drug} arm shows what that costs: \cref{tab:scale}
puts $199$ scored conditions behind it, drawn from $28$ treated wells and $40$
cell lines, so the smaller cluster count is the $28$ wells, the reference is
$t_{27}$ and not the normal, and the interval widens by $4.7\%$.

For statistics over pairs of conditions, such as the transfer coefficient,
neither grouping partitions the pairs, so those intervals use the
Fafchamps--Gubert dyadic estimator, under which two pairs are dependent when
they share either member.

Throughout, an estimate is quoted as value [lower, upper], a two-way
cluster-robust $95\%$ interval unless stated otherwise. DrEval~\cite{dreval}
identifies the same structure more broadly, citing Hurlbert~\cite{hurlbert}:
hundreds of thousands of cells resolving to a few hundred conditions over $50$
cell lines.

\section{The generation budget is part of the claim}

The second rule concerns generation validity. The budget is a claim-level
variable, not a runtime setting.

A $600$-token cap truncated $26\%$ of predictions in an early run, silently,
because nothing measured them. A prediction cut at the cap never reaches \DOWN{}
and is scored with its whole down-block missing, and no completion rate was
quoted beside the score.

\paragraph{What it cost} The optimal-transport arm of
\cref{sec:res-epsilon}, scored at that cap, gave a gap of
\ci{+0.0144}{-0.0051}{+0.0338} against its geometry-defined \texttt{orth}
partner, whose observed score is near chance. That interval spanned zero and the
arm was printed as inert.\corrmark{The printed negative is retracted in
\cref{sec:res-encoding}, in full and where it was made.} Re-run at the $1400$
tokens the residual arm had, it gives \ci{+0.0292}{+0.0144}{+0.0441}, and the
negative had to be withdrawn (\cref{sec:res-encoding}). The same cap had already
halved the residual arm's own effect. Wherever a paired comparison is scored on
variable-length predictions, then, the token budget is part of the claim, and a
null quoted without its budget cannot be read.

\paragraph{The rule that follows} Every generation-based evaluation computes the
\ENDCELL{}/\DOWN{} completion rate, the fraction of emitted symbols that are
valid panel genes, and the duplicate rate before any score is read from it, and
no score below comes from a run whose validity block was not inspected.

That is the apparatus. It settles what is measured, on what population, against
what reference and with what uncertainty, and reaches no verdict about the
model.
